% 【本文件写什么】一级组件 wateros-base，§ 级聚合
% - 组件职责与在内核中的位置：内核基础类型、地址表示与全局配置常量。
% - 聚合 crate os/components/wateros-base/src/lib.rs 的 pub mod 树与对外导出面
% - 子模块划分、与相邻组件的依赖边界，谁依赖谁、走哪条门面
% - 根 feature / 本组件 Cargo.toml feature 如何选用各 impl
% - 1～2 段综述后，由下方 \input 展开子模块；不在此逐条列举 api trait
% 【事实来源】
% - os/components/wateros-base/src/lib.rs、Cargo.toml
% - docs/exports/public-api/wateros-base.md
% - docs/exports/features/wateros-base.md
% - os/feature-tree.txt 中 wateros-base 相关段
% 【不写什么】
% - api-v0 契约细节 至 子目录 api/api-v0.tex
% - 各 impl 算法/平台细节 至 子目录 impl/*.tex

\section{基础类型组件}

\code{wateros-base}提供全内核共用的地址与标识薄类型，以及单核环境下的简单同步容器。它不决定堆有多大、定时器多长，只保证各处对物理地址、虚拟地址和hart编号的说法一致。配置数值单独放在同目录下的\code{base-config}子crate里，避免同一常量在两套源码里各写一遍。

根内核和其它组件通常同时依赖类型层与配置层：前者管怎么表示一个地址，后者管页多大、syscall最多几个参数、klog环多长这类 bring-up 常数。真机或自定义内存布局应以设备树探测结果为准，配置里的 QEMU 区间只是回退默认值。

\begin{rustcode}
// base/lib.rs：无 feature，四模块
pub mod addr;  // 物理/虚拟地址与页号包装
pub mod boot;  // 设备树物理基址别名
pub mod cpu;   // hart 标识
pub mod sync;  // 单核 RefCell 包装（多核须换锁）

// base-config：按子系统分文件
// syscall: MAX_SYSCALL_ARGS = 6
// mm: KERNEL_HEAP_SIZE、QEMU virt RAM/MMIO 回退区间
// task: 定时器周期、时间片上限
// fs/ipc/klog: 管道容量、页缓存规模、日志环大小
\end{rustcode}
